\documentclass[a4paper,12pt]{article}
\usepackage[finnish]{babel}
\usepackage[utf8]{inputenc}
\usepackage[T1]{fontenc}
\usepackage{lmodern}
\usepackage{eurosym}
\usepackage[pdftex,colorlinks]{hyperref}
\renewcommand{\thesubsection}{\arabic{subsection}}
\title{Etelä-Espoon sosialidemokraattinen työväenyhdistys r.y.}
\author{EED}
\date{21.01.2026}
\begin{document}
\maketitle
\tableofcontents
\section*{Johtokunnan kokous}
\subsection*{Paikka ja aika}
Tapiola Garden, keskiviikkona 21.01.2026 kello 18:00.
\section*{Läsnä}
Markku Sistonen, puheenjohtaja \\
Sinikka Suontio, varapuheenjohtaja \\
Arja Kivimäki, rahastonhoitaja \\
Martti Hellström, tiedotusvastaava \\
Jonne Hänninen \\
Matti Ikonen \\
Pauliina Sistonen \\
Mikko Nummelin, sihteeri
\section*{Pöytäkirja}
\subsection{Kokouksen avaus}
Puheenjohtaja avasi kokouksen kello .
\subsection{Kokouksen laillisuus ja päätösvaltaisuus}
Kokous todettiin lailliseksi ja päätösvaltaiseksi.
\subsection{Kokouksen ääntenlaskijoiden valinta}
Kokouksen ääntenlaskijoiksi valittiin.
\subsection{Johtokunnan järjestäytyminen}
Yhdistyksen sihteeriksi valittiin Mikko Nummelin. Yhdistyksen rahastonhoitajaksi valittiin . Yhdistyksen jäsenasiainhoitajaksi valittiin . Yrityksen tiedotusvastaavaksi valittiin .
\subsection{Falklammen uusi hinnasto}
Päätettiin uudesta Falklammen hinnastosta (alla):
\subsection{Jäsenasiat}
\subsection{Puoluekokousedustajat}
Yhdistys ehdottaa Helena Tiitistä ja Maria Guzeninaa esitettäväksi puoluekokousedustajiksi. Ylimääräinen jäsenkokous asiasta pidetään .
\subsection{Eduskuntavaaliehdokkaat}
Yhdistys ehdottaa Helena Tiitistä esitettäväksi eduskuntavaaliehdokkaaksi.
\subsection{Kevään tapahtumat}
\begin{itemize}
\item{Kreikan matka}
\end{itemize}
\subsection{Muut asiat}
\subsection{Kokouksen päättäminen}
Puheenjohtaja päätti kokouksen kello .
\end{document}
